\section{模型假设}

\begin{enumerate}
\item \textbf{电氢氨装置同步运行假设}：碱性电解槽（ALKEL）、质子交换膜电解槽（PEMEL）与合成氨装置作为整体同步运行，三者在运行时保持固定功率比例关系。该假设基于题目"电氢氨装置只有全额开机和停机两种运行方式"的表述，且制氢是为制氨服务，产业链上下游紧密耦合。

\item \textbf{产能线性扩容假设}：产能扩容过程中，仅电氢氨设备额定功率随产能规模呈线性同步提升，风光装机容量保持不变。72吨/日产能对应ALKEL 20 MW、PEMEL 20 MW、合成氨装置1.5 MW。该假设直接来源于题目"配套电氢氨装置的额定功率将随产能规模呈线性同步提升"的明确表述。

\item \textbf{功率损耗忽略假设}：不计园区内部功率损耗，即电能从发电侧到用电侧的传输效率为100\%。该假设由题目"不计园区功率损耗"明确给出。

\item \textbf{时段离散化假设}：以1小时为最小调度单元，每个时段内各设备功率恒定。该假设基于题目"园区运行分析时段为1小时"的规定，将连续时间问题离散化为24个决策时段。

\item \textbf{制氢产氢线性假设}：电解槽产氢速率与输入功率严格成线性关系，在功率连续可调范围（10\%--100\%额定功率）内效率恒定。该假设基于题目给出的额定功率与额定产氢速率的对应关系，且未提供效率-负荷曲线数据。
\end{enumerate}
